\section{Evaluating use, retention, and adaptation}
\label{sec:capability_eval}
We evaluate agent research capability across three dimensions:

\paragraph{Method selection and task quality.}
Tasks present distinct quantitative problems, requiring the agent to select tools,
set parameters, and interpret numerical outputs. We evaluate whether chosen algorithms
match input data properties (e.g., stationarity, regime type, asset universe).
Measuring method choice alongside returns separates sound reasoning from fortunate
price paths. Matched controls isolate the contributions of raw price tables, text
guidance, fixed tool outputs, and autonomous execution.

\paragraph{Context retention.}
Holding the base model, input data, and tool interfaces fixed, we compare full context,
recency truncation, and HiSTrim selective routing. We measure portfolio decision
quality, retention of critical price and instruction tokens, KV-cache memory footprint,
and generation latency. This determines whether pruning past interaction turns degrades
reasoning accuracy.

\paragraph{Experience and adaptation.}
Holding tool access and context budgets fixed, we evaluate sequential adaptation by
comparing unaugmented models, frozen-memory baselines, and our multi-timescale memory
module. All agents receive identical cleared returns. We track how tool choices shift
following losses, measuring turnover costs and performance across regime changes.
Factorial comparisons test interactions between context pruning and memory updates.

Evaluation metrics include forecast accuracy, portfolio Sharpe ratio, turnover, and
token cost. All model outputs and portfolio weights are saved and frozen prior to
scoring.
